% !TeX root = ../papers/nc-p3.tex
\section{二つの厚い色の方向が異なる窓の稠密部分}
\label{sec:klein-transverse-dense}
\begin{theorem}\label{thm:klein-transverse-dense}
二切断にそれぞれ階数2の色が二つ、階数1の色が二つあり、各切断の二つの階数2の色の差が異なるとする。
四つの階数1長方形が四行四列である場合、全正則窓は次の二変数表示とその位置の巡回・反転・文字変更で尽くされる。
和は $K$ の和とし、$b=i+j$, $e=j+k$, $l=i+e$ とする。厚い色を
$b\in\{0,1\}, e\in\{0,2\}$ に取る。総次数零の係数を通常1とし、
両条件を同時に満たす16語だけを
\[
 f=i_{\rm lo}+i_{\rm hi}+b_{\rm lo} e_{\rm hi}\pmod2,\qquad
 w_{ijkl}=\begin{cases}u&f=0,\\v&f=1\end{cases}
\]
に置き換える。正則性の必要十分条件は $u\ne v$ であり、$u=0$ または $v=0$ も許す。
最小周期は常に1である。反対階数は
$(u-1)(v-1)(uv+u+v-3)=0$ で12、その他で16である。
\end{theorem}

\begin{proof}[48成分の被覆]
階数1の全48成分の指数行列のSmith非零対角は35個の1であり、対角基底変更の階数は13である。
よってそれらを同時に1とできる。残る16語を辞書順に $x_0,\ldots,x_{15}$ とし、四つの厚い行列の
全三次小行列式の零点を取る。その極小素イデアルは48個である。
共通平行移動と位置別の $(t,t+h,t+h+k,t+k)$ の加算 ($h\in\{0,1\}$, $k\in\{0,2\}$)、
反転、巡回（奇数回のときは文字1,2も交換）で以下の五軌道になる。
\begin{center}
\begin{tabular}{c|l|c}
個数 & 代表の16係数の順序 & 始点0、四次の零文字次元\\\hline
16 & $(a,b,c,d,-1,\ldots,-1)$ & 11\\
16 & $(a,b,c,d,1,\ldots,1)$ & 11\\
8 & $(a,a,b,b,c,c,d,d,z,z,z,z,z,z,z,z)$ & 11\\
4 & $(a,a,a,a,b,b,b,b,c,c,c,c,d,d,d,d)$ & 13\\
4 & 定理の $(u,v)$ 表示 & 9（一般点）\end{tabular}
\end{center}
第三行には $-bcz+adz-a+b+c-d=0$ の一条件を課す。これは $a$ に関して原始的な一次式であり、
全ての代表は代数閉体への基礎拡大後も既約である。
\path{computations/verify_klein_transverse_dense.py} は極小素分解を五代表から独立に計算し、
各代表の軌道のイデアルとの一致をGröbner基底で検査する。初めの四代表について、
各関数体で文字次数零の関係行列の階数を計算したものが表の次元である。
標準Koszulの文字別Euler式の値は $2^2+2^2+1^2+1^2-1=9$ であり、特殊化で商次元は減らない。
従って最後の行だけが残る。
\end{proof}

\begin{proof}[二次完全交叉への変換]
各位置で $y_{a+2c}=e_a+(-1)^c e_{a+3}$ ($a,c=0,1$、$a+3$ は文字の和) による二文字ずつの
Fourier変換を施す。得たテンソルの文字総次数は零である。新しい文字の二コサイクル
$(-1)^{i_{\rm hi}j_{\rm lo}}$ で捻り、位置1,3で第3成分の符号を反転し、位置2で文字に1を加える。
得られるテンソル $S$ は完全対称である。全体を2で割った $S$ の縮約行列を
$S_b=(S_{i,i+b,k,k+b+1})_{ik}$ と書く。$S_0$ の第0,1列は等しく、第2,3列も等しい。
異なる二列の差と和は
\[
 4(1,1,-1,-1)^t,\qquad 2(u+1,v+1,u+1,v+1)^t
\]
である。$S_1=S_0^t$、$S_2$ の全列は $(u-1,v-1,u-1,v-1)^t$、
$S_3$ の全行は $(u-1,v-1,u-1,v-1)$ に等しい。これらの恒等式は全パラメータで成立する。

$u=v=1,-1$ では縮約階数が6に足りない。その他では $S$ の対称二次直交補空間は四次元である。
その共通零点 $[A:B:C:D]$ の四式は
\[
 A^2-C^2=B^2-D^2,\quad
 (v+1)(A^2+C^2)=(u+1)(B^2+D^2),\quad
 (v-1)AC=(u-1)BD,\quad AD=BC .
\]
最終式で $(A,B,C,D)=(sp,sq,tp,tq)$ と書く。第一式は
$(s^2-t^2)(p^2-q^2)=0$。$s^2=t^2$ なら中央二式の $p^2,q^2$ に関する係数行列の行列式は
符号を除き $2(u-v)$ である。$p^2=q^2$ なら $u\ne v$ のもとで
$s^2+t^2=st=0$ となって同じく解はない。$u=v$ なら $s=t=p=q=1$ が共通解である。
従って基点自由はちょうど $u\ne v$。定理\ref{thm:symmetric-frobenius-chart}と
二コサイクル捻りの左右Koszul複体の対応により、全点の正則性が決まる。
\end{proof}

\begin{proof}[周期と反対階数]
一位置巡回の後に1,2を交換し、四位置で $(0,1,1,0)$ を加えると、16語の条件も $f$ の値も元に戻る。
この四位置の基底変更を繰り返すと一移動と添字保持同型である。
反対切断の四色行列は行・列の置換を除き $J+\operatorname{diag}(u-1,v-1,v-1,u-1)$ であり、
行列式の零点は記載の三曲線である。$u\ne v$ のもとで二つの対角補正が零なら階数3、
全補正が非零なら対角行列に階数1を加えたものなので零行列式でも階数3である。
\end{proof}
